\documentclass[11pt]{article}

\usepackage[T1]{fontenc}

\usepackage[utf8]{inputenc}

\usepackage[margin=1in]{geometry}

\usepackage{graphicx}

\usepackage{amsmath}

\usepackage{amssymb}

\usepackage{booktabs}

\usepackage{array}

\usepackage{tabularx}

\usepackage{caption}

\usepackage{float}

\usepackage{microtype}

\usepackage{mathptmx}

\usepackage{natbib}

\usepackage[hidelinks]{hyperref}

\captionsetup{font=small,labelfont=bf}

\setlength{\parindent}{1.25em}

\setlength{\parskip}{0pt}

\title{Feasibility assessment for estimating smoking prevalence among Asian teens with the provided materials}

\author{}

\date{}

\begin{document}

\maketitle

\begin{center}\small\textit{Research Memo}\end{center}

\begin{abstract}

This memo assessed whether smoking prevalence among Asian teenagers could be estimated from the materials retrievable in this run. The workflow searched CDC-hosted youth tobacco sources and inspected candidate pages for machine-readable subgroup data, while also checking a user-provided GitHub hint. The accessible materials included a youth tobacco infographic landing page, an adult cigarette-smoking disparities page, and a GitHub repository hint that pointed to BRFSS 2015 diabetes files rather than smoking data. No retrievable dataset with teen-by-Asian subgroup smoking counts or denominators was obtained, so no original prevalence estimate could be computed without overclaiming. The run therefore documents feasibility constraints, source screening, and the reason a numeric result was not produced.

\end{abstract}

\section{Introduction}
The research question for this run was whether the prevalence of cigarette smoking among Asian teenagers could be estimated from the provided hints and directly retrievable public data sources. The intended focus was a race- and age-specific public health measure, but the availability of a valid estimate depended on finding a machine-readable dataset with both adolescent smoking outcomes and Asian subgroup structure. Because the task was constrained to the materials actually retrievable in this environment, the appropriate goal was a feasibility assessment rather than a substantive prevalence analysis. This memo records the search process, the sources checked, and the reason no estimate was produced.

\section{Methods}
I translated the theme into a concrete empirical question about smoking prevalence among Asian teenagers. I then searched CDC-hosted data sources for youth smoking datasets with race and ethnicity stratification, using queries focused on youth tobacco use, YRBS, NYTS, and Asian subgroup terms. I inspected the returned CDC pages to determine whether they exposed a machine-readable table suitable for original estimation in this run. I also checked the user-provided GitHub mirror hint and determined that it referred to a BRFSS 2015 diabetes dataset, which is not appropriate for estimating smoking prevalence among Asian teens. Because I could not retrieve a usable tabular dataset with teen-by-Asian subgroup smoking counts or denominators, I did not fabricate an estimate. I retained a retrieval log and a source-check table to document the feasibility assessment \ref{tab:artifact-1} \ref{tab:artifact-2}.

\section{Results}
No approved quantitative result was produced in this run. The source checks showed that the CDC youth tobacco search result was topic-relevant but did not yield a machine-readable teen-by-Asian subgroup dataset, while the other accessible CDC result was limited to adults and therefore did not match the study question \ref{tab:artifact-2}. The GitHub mirror hint was also not usable for this analysis because it pointed to BRFSS diabetes health indicators rather than smoking data for Asian teenagers \ref{tab:artifact-2}. The retrieval log documents these inspection steps and the resulting inability to compute a valid prevalence estimate from the materials obtained in this run \ref{tab:artifact-1}.

\begin{table}[tbp]
\centering
\caption{Run log documenting source searches, inspection steps, and why no valid smoking-prevalence estimate for Asian teens could be computed from the retrieved materials.}
\label{tab:artifact-1}
\begin{tabular}{>{\raggedright\arraybackslash}p{0.15\linewidth}}
\toprule
Attempted to identify a dataset for smoking prevalence among Asian teens using C \\
\midrule
Searches via web tool found adult smoking and youth tobacco infographic pages \\
Direct network retrieval inside python is disabled in this environment \\
\bottomrule
\end{tabular}
\end{table}

\begin{table}[tbp]
\centering
\caption{Table summarizing each candidate source checked in this run, its relevance, and why it was insufficient for original estimation.}
\label{tab:artifact-2}
\begin{tabular}{>{\raggedright\arraybackslash}p{0.15\linewidth}>{\raggedright\arraybackslash}p{0.15\linewidth}>{\raggedright\arraybackslash}p{0.15\linewidth}}
\toprule
source\_label & status & relevance \\
\midrule
CDC data search result: Youth And Tobacco Use Infographic & Found landing page only; no tabular subgroup data retrieved & topic-relevant but insufficient \\
CDC data search result: adult smoking infographic & Adult, not teen & insufficient for research question \\
GitHub BRFSS 2015 diabetes mirror & Available as hint but unrelated to smoking in Asian teens & not suitable primary data \\
\bottomrule
\end{tabular}
\end{table}

\section{Discussion}
The central finding of this memo is that the requested prevalence estimate could not be responsibly computed from the materials retrieved here. The search process identified relevant public-health sources, but relevance was not sufficient: the accessible CDC youth tobacco page did not expose the needed subgroup table, the adult smoking disparities page did not match the teen population, and the GitHub hint was unrelated to the outcome of interest \ref{tab:artifact-2}. In practical terms, the run shows that the analysis question is well-defined, but the available inputs were insufficient to support a valid original estimate. Reporting a numeric prevalence without a retrievable underlying dataset would have risked overstatement. The value of the run is therefore procedural: it clarifies what was searched, what was found, and why the work stopped short of quantitative estimation.

\section{Limitations}
No machine-readable dataset containing smoking outcomes for teenagers stratified specifically for Asian youth was successfully retrieved in this run. The CDC search results that were accessible included an infographic landing page and adult smoking material, but not a directly usable teen subgroup table for original estimation. The GitHub mirror hint available in the prompt pointed to a BRFSS 2015 diabetes dataset, which is not appropriate for estimating smoking prevalence among Asian teens. Python in this environment could not directly fetch external URLs, so analysis depended on the web retrieval tool; without a downloadable tabular file exposed there, computation was not possible. Because meaningful original estimation was not possible with the materials actually retrieved, returning a numeric prevalence would risk overclaiming.

\bibliographystyle{plainnat}
\bibliography{references}

\end{document}
